\documentclass[a4paper, amsfonts, amssymb, amsmath, onecolumn, showkeys, nofootinbib, twoside, superscriptaddress]{revtex4-2}
\usepackage[english]{babel}
\usepackage[utf8]{inputenc}
\usepackage[colorinlistoftodos, color=green!40, prependcaption]{todonotes}
\usepackage{cancel,amsmath}
\usepackage{amsthm}
\usepackage{mathtools}
\usepackage{physics}
\usepackage{xcolor}
\usepackage{graphicx}
\usepackage{float}
\usepackage{soul}

\usepackage{gensymb}
\usepackage{multirow}
\usepackage[left=23mm,right=13mm,top=35mm,columnsep=15pt]{geometry} 
\usepackage{adjustbox}
\usepackage{placeins}
\usepackage[T1]{fontenc}

\usepackage{csquotes}

\usepackage[pdftex, pdftitle={Article}, pdfauthor={Author}]{hyperref} % For hyperlinks in the PDF

\renewcommand\floatpagefraction{0.8} %% default value: 0.5
\renewcommand\topfraction{0.8}       %% default value: 0.7
\bibliographystyle{apsrev4-1}

% \usepackage{changes}
% \usepackage[final]{changes}

% cleveref! \cref{x} gives the result of "ref type"~\ref{x}
\usepackage[capitalize]{cleveref}
\DeclareMathOperator*{\argmax}{arg\,max}


\begin{document}
\title{First-Passage Approach to Optimizing Perturbations for Improved Training of Machine Learning Models}


\author{Sagi Meir}
    \affiliation{School of Chemistry, Tel Aviv University, Tel Aviv 6997801, Israel.}
    \affiliation{The Center for Physics and Chemistry of Living Systems, Tel Aviv University, Tel Aviv 6997801, Israel.}
\author{Tommer D. Keidar}
    \affiliation{School of Chemistry, Tel Aviv University, Tel Aviv 6997801, Israel.}
    \affiliation{The Center for Physics and Chemistry of Living Systems, Tel Aviv University, Tel Aviv 6997801, Israel.}
\author{Shlomi Reuveni}
    \affiliation{School of Chemistry, Tel Aviv University, Tel Aviv 6997801, Israel.}
    \affiliation{The Center for Physics and Chemistry of Living Systems, Tel Aviv University, Tel Aviv 6997801, Israel.}
    \affiliation{The Center for Computational Molecular and Materials Science, Tel Aviv University, Tel Aviv 6997801, Israel.}
\author{Barak Hirshberg}
    \email{hirshb@tauex.tau.ac.il}
    \affiliation{School of Chemistry, Tel Aviv University, Tel Aviv 6997801, Israel.}
    \affiliation{The Center for Physics and Chemistry of Living Systems, Tel Aviv University, Tel Aviv 6997801, Israel.}
    \affiliation{The Center for Computational Molecular and Materials Science, Tel Aviv University, Tel Aviv 6997801, Israel.}


\begin{abstract}
    Machine learning models have become indispensable tools in applications across the physical sciences. Their training is often time-consuming, vastly exceeding the inference timescales. Several protocols have been developed to perturb the learning process and improve the training, such as shrink and perturb, warm restarts, and stochastic resetting. For classifiers, these perturbations have been shown to result in enhanced speedups or improved generalization. However, the design of such perturbations is usually done \textit{ad hoc} by intuition and trial and error. To rationally optimize training protocols, we frame them as first-passage processes and consider their response to perturbations. 
    We show that if the unperturbed learning process reaches a quasi-steady state, the response at a single perturbation frequency can predict the behavior at a wide range of frequencies. 
    We employ this approach to a CIFAR-10 classifier using the ResNet-18 model and  identify a useful perturbation and frequency among several possibilities. 
    Our work allows optimization of perturbations for improving the training of machine learning models using a first-passage approach.
\end{abstract}

\keywords{machine learning, neural networks, perturbations, stochastic resetting, first-passage, response theory}

\maketitle


\section{Introduction}
Machine learning algorithms and deep neural network (NN) models have become powerful tools across chemistry and physics to tackle challenges that were once computationally or experimentally prohibitive. Examples include protein structure prediction~\cite{jumper2021highly}, entropy calculation for physical systems~\cite{nir2020machine}, machine learning potentials~\cite{behler2007generalized}, and many others~\cite{carleo2017solving, noe2019boltzmann, mate2024neural, anelli2024robust, lagemann2021deep, ravuri2021skilful, wang2020machine, rupp2012fast, tsai2020learning, geiger2013neural}. Unfortunately, these capabilities come with a time-consuming price spent on model training. Second-order optimization algorithms are unfeasible for deep NN~\cite{fukumizu2000local,loshchilov2016sgdr}, and currently, stochastic gradient descent (SGD) and its variants, remain a cornerstone of NN training~\cite{peng2019accelerating, loshchilov2016sgdr, nguyen2019machine, rumelhart1986learning, duchi2011adaptive, sutskever2013importance, zeiler2012adadelta, kingma2014adam}. 

SGD by itself can be slow and prone to getting stuck in local minima. Several common practices are used to address these limitations. First is the incorporation of momentum and adaptive learning rates as done by modern optimizers, usually leading to faster convergence and improved performance~\cite{sutskever2013importance, zeiler2012adadelta, kingma2014adam, leimkuhler2019partitioned}. The second, which is the focus of this work, is to ``perturb'' the training process, which may lead to speedup or improved generalization. Examples include shrink \& perturb (S\&P)~\cite{ash2020warm, zaidi2023does}, warm restarts~\cite{loshchilov2016sgdr}, stochastic resetting (SR)~\cite{bae2024stochastic} and continual backpropagation~\cite{dohare2024loss}. 

Ash et al.~\cite{ash2020warm} first suggested the S\&P protocol to improve generalization during online learning. More recently, Zaidi et al.~\cite{zaidi2023does} investigated when such re-initialization might help the training, and concluded that while it is clear S\&P is helpful in some cases, a general theory of why it works is missing. Bae et al.~\cite{bae2024stochastic} applied SR  to a dynamically-updated checkpoint during the training of an NN with a resetting time sampled from an exponential distribution. They concluded that their strategy decreases overfitting on noisy labels and leads to better generalization. Yet, predicting the optimal resetting rate a priori was not possible. Loshchilov et al.~\cite{loshchilov2016sgdr} showed that a protocol with warm restarts converges up to four times faster to the same test accuracy compared to standard SGD training. 
Overall, despite their usefulness, a theory capable of predicting the effect of such perturbation protocols is missing, and their design is often done by intuition and empirical trial and error. 

The resemblance of SGD to Langevin dynamics~\cite{feng2021inverse,cheng2020stochastic} has been utilized to describe and analyze the learning process of NNs with concepts and methods from thermodynamics and statistical physics~\cite{seung1992statistical, choromanska2015loss, alemi2018therml, stephan2017stochastic, feng2021inverse, jules2023charting, zdeborova2020understanding, karniadakis2021physics, carleo2019machine}.
For example, Feng et al. found an inverse fluctuation-dissipation relation between weights' fluctuations and the flatness of the loss landscape. Based on their finding, they developed an algorithm that delays catastrophic forgetting in sequential learning tasks~\cite{feng2021inverse}.

Another example is the work of Stephan et al. where they showed that SGD with a constant learning rate can be used as an approximate Bayesian posterior inference algorithm~\cite{stephan2017stochastic}. Their result was obtained by viewing SGD as a Markov chain with a stationary distribution.
However, despite the advancements in applying statistical mechanical tools in machine learning, they have not been previously used to design perturbations to improve the performance of SGD optimizers.

Recently, Keidar et al. developed a response theory to predict how rare perturbations affect the completion of an arbitrary stochastic process~\cite{keidar2024universal}. Their theory focused on first-passage processes~\cite{Redner_2001,metzler2014first,bray2013persistence,reuveni2016optimal, pal2017first, evans2020stochastic, blumer2022stochastic, blumer2024combining}, in which a stochastic process has a predefined distribution of initial conditions and a well-defined target. The stochastic nature of the process gives rise to a distribution of first-passage times, i.e., the first time the process reaches its target. In this work, we treat the SGD-based training of NNs up to a predefined target accuracy as a first-passage process. We treat protocols such as S\&P, warm restarts, and stochastic resetting as perturbations to the first-passage process that occur every $P$ epochs during the training. We employ the theory of Keidar et al. for the first time to analyze and improve the training of NN models using such perturbations. We show that, given a set of protocols, we can determine which would lead to the highest acceleration and identify the optimal perturbation time interval $P$. We focus on three types of protocols: S\&P, partial re-initialization of small weights (partial SR), and full SR. Finally, we suggest a methodology for testing and analyzing new perturbation protocols.


\section{Theory}

In this work, we treat the NN training in the absence of a perturbation as a stochastic process that is characterized by a propagator $G(\boldsymbol{\theta},t)$, representing the probability of being in state $\boldsymbol{\theta}$ at time $t$. We emphasize that $\boldsymbol{\theta}$ characterizes the overall state of the system, i.e., it may represent the weights, biases, hyperparameters, etc. We focus on first-passage processes (see \cref{fig:LR_Theory}A), and define the first-passage time (FPT)  as the first instance in which the test accuracy reaches a certain threshold, which is treated as an absorbing boundary condition. Namely, once the test accuracy of a model reaches the target, training is stopped. Since it is a stochastic process, there will be an FPT distribution. We denote with $T$ the random variable representing the FPT of the unperturbed training process. We define the survival probability $\Psi_T(t)$ as the fraction of models which did not reach the threshold at time $t$,
\begin{equation}
    \Psi_T(t) = \int_\Theta G(\boldsymbol{\theta},t)\mathrm{d}\boldsymbol{\theta},
\label{eq:Psi-definition}
\end{equation}
where $\Theta$ symbolizes that the integration domain is over all possible states. 


\begin{figure}[h]
\begin{center}
\centerline{\includegraphics[width=0.7\linewidth]{LR_theory_pic.pdf}}
\caption{Training NNs as a first-passage process. Panel A presents the test accuracy as a function of the number of epochs for training without perturbations (orange line). Here, $T$ is the FPT to reach the target accuracy (dashed blue line). Panel B presents the test accuracy as a function of the number of epochs for training with a perturbation every $P$ epochs (orange line). $T_P$ is the perturbed FPT to the target accuracy (dashed blue line) and $\tau_P(\boldsymbol{\theta})$ is the residual time to reach the target accuracy after the first perturbation.}
\label{fig:LR_Theory}
\end{center}
\end{figure}

\cref{fig:LR_Theory}B shows the training process with a perturbation that is applied every $P$ epochs. We denote by $T_P$ the random variable representing the FPT of the perturbed process to the target accuracy. The perturbation can be of any form, e.g., it can affect the network weights, hyperparameters, activation functions, or any other component of the training. Keidar et al. showed that the perturbed FPT is connected to the unperturbed FPT through~\cite{keidar2024universal},
\begin{equation}
T_P= \begin{cases}T & \text { if } T \leq P, \\ P+\tau_P(\boldsymbol{\theta}) & \text { if } T>P,\end{cases}
\label{eq:T_P_cases_description}
\end{equation}
where $\tau_P(\boldsymbol{\theta})$ is the random variable representing the residual number of epochs it takes to reach the target after the perturbation has been first applied (see \cref{fig:LR_Theory}B). It depends on $\boldsymbol{\theta}$, the state of the network at time $P$, which might be different for every realization of the process. 

Applying the law of total expectation to \cref{eq:T_P_cases_description} (see the derivation in the appendix), we obtain the mean FPT under perturbations, $\mathbb{E}[T_P]$,
\begin{equation}
    \mathbb{E}[T_P]=\sum_{t=0}^{P-1}\Psi_T(t)  + \Psi_T(P) \Bar{\tau}_{P},
\label{eq:LR-prediction}
\end{equation}
where $\Bar{\tau}_{P}$ is given by
\begin{equation}
    \Bar{\tau}_{P} = \int_{\Theta} \mathbb{E}[\tau_P(\boldsymbol{\theta})] \frac{G(\boldsymbol{\theta},P)}{\Psi_T(P)} \mathrm{d}\boldsymbol{\theta}.
\label{eq:tau_P-definition}
\end{equation}
In \cref{eq:tau_P-definition}, $\mathbb{E}[\tau_P(\boldsymbol{\theta})]$ is the average residual time after the first perturbation over all possible noise realizations of the stochastic training process after the perturbation was applied to a given state $\boldsymbol{\theta}$ at time $P$.  Hence, $\Bar{\tau}_{P}$ can be understood as the average of $\mathbb{E}[\tau_P(\boldsymbol{\theta})]$ over all possible $\boldsymbol{\theta}$ generated by the unperturbed process at time $P$ right before the perturbation has been applied.

\cref{eq:LR-prediction} formally decomposes the mean FPT with the perturbation to two contributions. The first term, $\sum_{t=0}^{P-1}\Psi_T(t)$, sets a lower bound on $\mathbb{E}[T_P]$ that only depends on the unperturbed process. The second term encodes all the effects of the perturbation on the mean FPT. These equations allow, in principle to predict the effect of a perturbation on the training for all $P$. 

Consider, for example, the specific case of SR as the perturbation. Then, every $P$ epochs, the state of the system is restarted by resampling the initial conditions for training. As a result, for any training process, $\tau_P(\boldsymbol{\theta})$ does not depend on the state at time $P$, and is simply an independent and identically distributed copy of $T_P$. In that case, $\Bar{\tau}_P=\mathbb{E}[T_P]$, and substituting it into \cref{eq:LR-prediction} gives~\cite{eliazar2020mean},
\begin{equation}
    \mathbb{E}[T_P]_{SR}=\frac{1}{1-\Psi_T(P)}\sum_{t=0}^{P-1}\Psi_T(t) .
\label{eq:SR-prediction}
\end{equation}
\cref{eq:SR-prediction} shows that, for SR, the mean FPT with the perturbation can be predicted entirely from the survival probability of the unperturbed learning process. 
To derive \cref{eq:SR-prediction}, we used the specific properties of SR. For other perturbations, one must assume something about the dynamics of the underlying training process as we do in the next section. 


\section{Results and Discussion}

\subsection{The quasi-steady state hypothesis}
\label{sec:qss-theory}
To proceed beyond \cref{eq:LR-prediction}, we consider a stochastic process in which the propagator reaches a quasi-steady state (QSS) after a typical relaxation time $t_r$,
\begin{equation}
    G(\boldsymbol{\theta},t)\simeq \phi(\boldsymbol{\theta})\Psi_T(t) \quad \text { for } t>t_r.
\label{eq:quasi-steady state}
\end{equation}
In \cref{eq:quasi-steady state}, $\phi(\boldsymbol{\theta})$ is a time-independent probability density function of the network state $\boldsymbol{\theta}$. In other words, a QSS is defined by constant relative populations of all states $\boldsymbol{\theta}$, although the overall number of models that survived declines over time through $\Psi_T(t)$~\cite{nitzan2024chemical}, as illustrated in \cref{fig:qss_drawing}. The QSS approximation is commonly applied to analyze chemical kinetics~\cite{nitzan2024chemical, ji2021stiff, kramers1940brownian, hanggi1986escape, hanggi1990reaction}, but has not been used to explain NN training before. Next, we use it to show that for rare perturbations, the mean residual time $\Bar{\tau}_P$ is independent of $P$. Therefore, by using samples of the training at a single value of $P$, i.e., a single perturbation frequency, we can predict the mean FPT at a wide range of perturbation frequencies.

\begin{figure}[h]
\begin{center}
\centerline{\includegraphics[width=0.69\linewidth]{QSS_drawing_2ndVER.pdf}}
\caption{An illustration of a QSS of a propagator $G(\boldsymbol{\theta},t)$ at times $t_2>t_1>t_r$. At QSS, the probability density function $\phi(\boldsymbol{\theta})$ (purple lines) maintains the same shape, while the fraction of models that survived (blue areas) declines over time through panels A-B, i.e., $\Psi_T(t_1)>\Psi_T(t_2)$.}
\label{fig:qss_drawing}
\end{center}
\end{figure}


To show this, we consider three timescales in the system: the relaxation time $t_r$, the perturbation time interval $P$ and the mean residual time $\Bar{\tau}_P$. Considering a perturbation that is done at a time $P > t_r$, and utilizing QSS, we observe that the perturbation acts on a state sampled from $\phi(\boldsymbol{\theta})$ irrespective of $P$. We can then plug \cref{eq:quasi-steady state} into \cref{eq:tau_P-definition} and get
\begin{equation}
    \Bar{\tau}_{P} = \int_{\Theta} \mathbb{E}[\tau_P(\boldsymbol{\theta})] \phi(\boldsymbol{\theta}) \mathrm{d}\boldsymbol{\theta}.
\label{eq:tau-indep-P_step1}
\end{equation}
Note that the mean residual time for completion after the perturbation is applied still depends on $P$ through $\mathbb{E}[\tau_P(\boldsymbol{\theta})]$. Next, if the perturbation is rare enough, i.e., $P > \Bar{\tau}_P$, we effectively perturb the system only once on average before the training is complete. As a result, $\tau_P(\boldsymbol{\theta})=\tau(\boldsymbol{\theta})$, i.e., does not depend on $P$, leading to
\begin{equation}
    \Bar{\tau}_{P} \approx \Bar{\tau} = \int_{\Theta} \mathbb{E}[\tau(\boldsymbol{\theta})] \phi(\boldsymbol{\theta}) \mathrm{d}\boldsymbol{\theta} \quad \text { for } P>\max\left(t_r, \Bar{\tau}_P\right).
\label{eq:tau-indep-P}
\end{equation}
%\textcolor{red}{Note that we obtained \cref{eq:tau-indep-P} by assuming a QSS, but hypothetically there could be non-QSS scenarios for which \cref{eq:tau-indep-P} would hold. Fortunately, our method applies also to such cases.} 
Plugging \cref{eq:tau-indep-P} into \cref{eq:LR-prediction} leads to
\begin{equation}
    \mathbb{E}[T_P]=\sum_{t=0}^{P-1}\Psi_T(t)  + \Psi_T(P) \Bar{\tau}.
\label{eq:LR-prediction-rarepert}
\end{equation}
\cref{eq:LR-prediction-rarepert} sets the recipe for designing useful perturbations and the frequencies at which to apply them. It tells us to perform simulations applying the perturbation only once at some $P^*>\max\left(t_r, \Bar{\tau}_P\right)$. These simulations provide the unbiased survival function until time $P^*$ and $\Bar{\tau}$. Combining the two, we can predict the mean FPT for all values $\max\left(t_r, \Bar{\tau}_P\right) \leq P \leq P^*$. 

Below, we first test the QSS hypothesis when training a CIFAR-10~\cite{krizhevsky2009learning} classifier using the ResNet-18 model~\cite{he2016deep}. Then, for two perturbations, S\&P and partial SR, we show that indeed $\Bar{\tau}_{P} \approx \Bar{\tau}$ for a wide range of $P$. Finally, we use \cref{eq:LR-prediction-rarepert} to predict the mean FPT under these perturbations at a wide range of $P$, and benchmark them against brute-force training. We show that this procedure can be used to select a perturbation and time interval $P$ that lead to faster training.


\subsection{Experimental test of the quasi-steady-state}
A naive approach to demonstrate a QSS when training an NN model would be to sample several learning trajectories while keeping track of the parameters to obtain $G(\boldsymbol{\theta},t)$. However, it is impractical to keep track of millions of parameters, and the finite sample size introduces noise to the estimation of the distribution of $G(\boldsymbol{\theta},t)$. This noise increases with the dimensionality of $\boldsymbol{\theta}$, making it infeasible to demonstrate a QSS convergence from the trajectory data of $\boldsymbol{\theta}$. 
Instead, we draw inspiration from enhanced sampling of free energy surfaces~\cite{barducci2011metadynamics, valsson2014variational, bussi2020using} and define a collective variable,  $A(\boldsymbol{\theta})$, representing the state of the system and look at its marginal distribution, 
\begin{equation}
    G(A,t) = \int_{\Theta} G(\boldsymbol{\theta},t)\delta(A(\boldsymbol{\theta})-A)\mathrm{d}\boldsymbol{\theta}.
\label{eq:acc_propagator-def}
\end{equation}
In practice, we will use the test set accuracy as a collective variable. 
If $G(\boldsymbol{\theta},t)$ reaches a QSS, and we plug \cref{eq:quasi-steady state} into \cref{eq:acc_propagator-def}, we get a QSS also in $G(A,t)$,
\begin{equation}
    G(A,t) \simeq \Psi_T(t)\int_{\Theta} \phi(\boldsymbol{\theta})\delta(A(\boldsymbol{\theta})-A)\mathrm{d}\boldsymbol{\theta} = \Psi_T(t)\phi(A), \quad \text { for } t>t_r.
\label{eq:acc_quasi-steay}
\end{equation}
In \cref{eq:acc_quasi-steay}, $\phi(A)=\int_{\Theta} \phi(\boldsymbol{\theta})\delta(A(\boldsymbol{\theta})-A)\mathrm{d}\boldsymbol{\theta}$ is the time-independent probability density function of $A$. Therefore, we will use the dynamics of $G(A,t)$ as a proxy to justify the QSS hypothesis.


We used the ResNet-18 architecture as our NN model for the classification task of the CIFAR-10 dataset. We trained 1000 models using random initialization and plotted the test accuracy as a function of the number of epochs, which we will refer to as trajectories. We used standard SGD training, without perturbations. See the Computational Details section for the full training setup. To obtain the propagator $G(A,t)$ from the trajectory data for a specific target test accuracy, we treat it as an absorbing boundary. In other words, trajectories that reached the target accuracy are stopped and excluded from $G(A,t)$ at later times. The survival $\Psi_T(t)$ is estimated by the fraction of models that did not reach the target at epoch $t$ of the training (see \cref{eq:Psi-definition}). 
%For example, if the target test accuracy is set to an unreachable value then $\Psi_T(t)=1,\, \forall t$. Particularly in this case, any QSS is a steady state.

In \cref{fig:acc_densities_over_time}A, we start by analyzing the specific case of a target test accuracy of 75\% (dashed blue line). For this accuracy, no model reaches the target up to 100 epochs, i.e., $\Psi_T(t)=1 \, \forall t$ (blue line \cref{fig:acc_densities_over_time}C). The violins of \cref{fig:acc_densities_over_time}A represent the density distribution of $A$ as a function of time. It narrows down with time, until epoch $\sim20$, and then remains with a fixed shape and height, which suggests that $G(A,t)=\phi(A)$ from that epoch onwards. 
%We define the relaxation time according to the convergence of the mean, i.e., $t_r\sim20$ epochs. 
Besides this qualitative shape comparison, we also quantitatively compare the cumulative distribution functions (CDFs) of $A$ to the average CDF of $A$ over epochs 20 to 100 with the Kolmogorov–Smirnov (KS) test~\cite{massey1951kolmogorov, tiwary2015role, blumer2024short}, shown in \cref{fig:acc_densities_over_time}D (blue line). We define the relaxation time as having a p-value larger than 0.05 in the KS test, which gives a relaxation time of $t_r=17$ epochs. Alternative statistical tests result in similar relaxation times (see the appendix).

\begin{figure}[h]
\centering
    \begin{minipage}[t]{0.78\linewidth}
    \centering
    \includegraphics[width=\linewidth]{acc_densities_over_time.pdf}
    \end{minipage}
    \hspace{-0.0262\linewidth}
    \begin{minipage}[t]{0.2195\linewidth}
    \centering
    \includegraphics[width=\linewidth]{qss_distance.pdf}
    \end{minipage}
\caption{Experimental evidence for a quasi-steady-state. Panels A and B are violin plots of the density distributions of the test accuracy as a function of the number of epochs, for models that did not reach the target. Blue dots represent the mean accuracy and the blue areas are the distributions rotated by $90^{\circ}$. The dashed blue lines are the target accuracies, set to 75\% and 72\% in panels A and B, respectively. In both cases, after $\sim$20-30 epochs, the system reaches a QSS, i.e., $G(A,t)\simeq \phi(A)\Psi_T(t)$. 
Panel A is a special case in which $\Psi_T(t)=1$, while in panel B, $\Psi_T(t)$ decays slowly. 
Panel C shows the survivals $\Psi_T(t)$ for the 75\% (blue), 72\% (orange), and 71.5\% (green) target test accuracies. 
Panel D is the Kolmogorov–Smirnov (KS) test between the CDFs of the accuracy $A$ at different epochs to the average CDF of $A$ over epochs 20 to 100. 
}
\label{fig:acc_densities_over_time}
\end{figure}


In \cref{fig:acc_densities_over_time}B we set the target test accuracy to 72\% which leads to a decay of $\Psi_T(t)$ with the number of epochs, i.e., some trajectories reach the target during the training (orange line in \cref{fig:acc_densities_over_time}C). Although there is a decline in the number of models that did not reach the target, after $\sim30$ epochs the shape of the density distribution again remains approximately fixed while the height decreases, as expected from a QSS (\cref{eq:acc_quasi-steay}). 
%Additional, indication to a QSS is the orange line in \cref{fig:acc_densities_over_time}C that changes its decay behavior at $t_r$ to a slow decay. 
We compare the CDFs of $A$ to the average CDF of $A$ over epochs 20 to 100 with the KS test~\cite{massey1951kolmogorov, tiwary2015role, blumer2024short}, conditioned on model survival.
Similarly to the previous case, we plot the p-value of the KS-test in \cref{fig:acc_densities_over_time}D (orange line) and obtain $t_r = 28$ epochs.
We also checked the QSS hypothesis for a lower target accuracy of $71.5\%$ and plot the survival function and KS-test p-value in \cref{fig:acc_densities_over_time}C,D (green lines), respectively. Again, we obtain a similar relaxation time of $t_r = 26$ epochs.
All these examples are consistent with our hypothesis that the training of a CIFAR-10 classifier using a ResNet-18 model reaches a QSS for various target accuracies.


\subsection{Experimental test that the residual time is $P$-independent}
We now show that if $P$ is large enough, introducing a perturbation would result in $P$-independent residual times, according to \cref{eq:tau-indep-P}. Here, we show it experimentally by focusing on the 72\% target test accuracy  case (\cref{fig:acc_densities_over_time}B) where the system reaches a QSS after $t_r=28$ epochs. We analyze the dependence of the residual time to reach the target on the perturbation time $P$. To do so, we take only the models that did not reach the target after $P=100,\,50,\,20,$ and 10 epochs, introduce the perturbation, and continue the training until reaching the target. We consider two perturbations, S\&P and partial SR. 

The S\&P perturbation takes the weights and biases of the NN at the moment it is applied, shrinks them by a factor $\lambda$ and adds an independent and identically distributed copy sampled from the initial conditions distribution, multiplied by a factor $\gamma$. We used the same values as in ref. \cite{zaidi2023does}, $\lambda=0.4$ and $\gamma=0.1$. The partial SR perturbation takes only a fraction of the parameters of the NN and re-initializes them to an independent and identically distributed copy sampled from the initial conditions distribution. We chose to initialize 30\% of the weights, those with the smallest absolute value at the moment the perturbation is applied.

In \cref{fig:s_and_p_constant_tau}, we first plot the mean test accuracy of the models as a function of the epoch. The top and bottom rows correspond to the S\&P and the partial resetting perturbations, respectively. In all cases, we find that the residual time of the mean trajectory to reach the target is smaller than $P$ and that the perturbation is only applied once before reaching the target. 
As a result, when $P>t_r$, we expect that the mean residual time will be $P$-independent, as explained in \cref{sec:qss-theory}. 
 In panels E and J we center all the mean trajectories with respect to their perturbation time $P$. We find that all curves approximately collapse on top of each other regardless of $P$, reaching the target accuracy after the same number of epochs since the perturbation has been applied (gray area). We emphasize that the residual time is different for the two perturbations, as shown in panels E and J, but is $P$-independent in both cases.

\begin{figure}[h]
\begin{center}
\centerline{\includegraphics[width=0.99\linewidth]{constant_tau.pdf}}
\caption{The mean trajectory of the test accuracy and its top and bottom deciles (shaded areas), for models that did not reach the target after $P$ epochs. From left to right $P=100$, $50$, $20$, $10$ (see legend). The dashed black lines represent 72\% test accuracy (the target). At epoch number $P$ we apply the S\&P protocol (upper row, panels A-E) or the partial SR protocol (bottom row, panels F-J). Panels E and J center the mean trajectories with respect to their $P$. The gray areas are from the moment of perturbation until a the first time that the average trajectory reaches the target accuracy.}
\label{fig:s_and_p_constant_tau}
\end{center}
\end{figure}

Next, instead of plotting the average training curve, we plot the value of the residual time $\Bar{\tau}_P$ as a function of $P$ in \cref{fig:tau_P_vs_P} for a larger range of $P=1\text{--}100$. 
We observe that $\Bar{\tau}_P$ is roughly constant over nearly two orders of magnitude in $P$ for both perturbations. Surprisingly, $\Bar{\tau}_P \approx \Bar{\tau}$ even for values of $P$ that are smaller than the relaxation time, beyond the expected range according to \cref{eq:tau-indep-P}. However, when $P$ becomes smaller than the residual time ($\sim3$ and $\sim5$ for S\&P and partial SR, respectively), denoted by the shaded area in \cref{fig:tau_P_vs_P}, the process is perturbed more than once on average before reaching the target and the residual time is no longer $P$-independent. 


\begin{figure}[h]
\begin{center}
\centerline{\includegraphics[width=0.4\linewidth]{tau_P_vs_P.pdf}}
\caption{The mean residual time $\Bar{\tau}_P$ versus the perturbation time. The values of $\Bar{\tau}_P$ for S\&P (orange circles) and partial SR (green triangles) are approximately constant when $P>\Bar{\tau}_P$ (outside the yellow area). Dashed dotted lines are averages of all $\Bar{\tau}_P<P$.}
\label{fig:tau_P_vs_P}
\end{center}
\end{figure}



\subsection{Prediction of the mean FPT}
In this section, we show how to predict the mean FPT with a perturbation for a wide range of $P$ from as few experiments as possible. 
We begin by assuming that several training trajectories are already available, and we wish to assess whether applying a perturbation is worthwhile, and select at what time interval to apply it. The simplest perturbation to analyze is SR. It comes at no additional cost, through \cref{eq:SR-prediction}, since we already sampled the survival function. To quantify the acceleration in reaching a target accuracy by applying a perturbation, we define the speedup as the ratio between the mean FPT without and with the perturbation, respectively. We plot the predicted speedup for SR in \cref{fig:Speedup} (blue line) and compare it with numerical experiments (blue squares). We find that SR leads to a maximal speedup of $\sim3$ using a perturbation time interval of $\sim20$. 

To go beyond SR to other perturbations that could potentially lead to higher speedups, we must measure their residual time.
A naive approach of simply evaluating $\Bar{\tau}_P$ for every $P$ and using \cref{eq:LR-prediction} would mean running the training for all perturbations at all $P$ values, which is costly. Instead, we would like to use the fact that the residual time is $P$-independent. To that end, we sample the training with the perturbation applied once at some $P^*$ that is large enough. We then use \cref{eq:LR-prediction-rarepert} to predict the mean FPT for a wide range of $P < P^*$. 
We plot the predicted speedup in \cref{fig:Speedup} for the S\&P (orange line) and partial SR (green line) perturbations, using their value of $\Bar{\tau}_{P^*}$ at $P^*=100$. Going from right to left on the $P$ axis, the predictions fit the experimental values of the S\&P (orange circles) and partial SR (green triangles) protocols for $P\geq\Bar{\tau}_{P^*}$, as expected.

Using these predictions, we identify that S\&P leads to a speedup of $\sim16$ at a perturbation time interval of 3, while partial SR leads to a speedup of $\sim8$ at $P=7$. These predictions correctly identify S\&P as the preferred strategy that leads to the highest speedup. Although the experimental maximal speedup is slightly larger ($\sim21$), our predictions are a computationally efficient way of determining useful perturbations and time intervals that lead to high speedups.


\begin{figure}[h]
\begin{center}
\centerline{\includegraphics[width=0.4\linewidth]{Speedup_noMomentum.pdf}}
\caption{The speedup gained by using different perturbation protocols versus the perturbation time. Orange and green solid lines are the theoretical predictions of \cref{eq:LR-prediction-rarepert} for S\&P and partial SR, respectively. For the predictions, we used the value of $\Bar{\tau}_{P^*}$ for $P^*=100$. We plot the predictions only down to $P=\Bar{\tau}_{P^*}$.
The blue solid line is the theoretical prediction for the SR case (\cref{eq:SR-prediction}). Symbols represent brute force training at every value of $P$.}
\label{fig:Speedup}
\end{center}
\end{figure}

We suggest a general methodology that can examine different kinds of perturbations and pick the optimal one and an efficient $P$. Consider for example a classification task for which the architecture of the NN and the hyperparameters of the SGD-based optimization are given. Then, it is common practice to train in parallel an ensemble of models and terminate any training process that did not reach a minimal desired test accuracy after a tolerance time of $P^*$ epochs. 
After this process is completed, we propose to estimate the survival function, $\Psi_T(t)$ for $t\leq P^*$. This is already enough information to evaluate whether SR leads to any acceleration. 
For other perturbations, evaluate first if the training accuracy of the unperturbed process  reaches a QSS and determine $t_r < P^*$ using the KS test. If that is the case, introduce several candidate perturbations to the models that did not reach the target accuracy and measure $\Bar{\tau}_{P^*}$ for each one. 
This way, $\Psi_T(t\leq P^*)$ and $\Bar{\tau}_{P^*}$ are obtained simultaneously for several perturbation protocols. According to \cref{eq:LR-prediction-rarepert}, the perturbation with the minimal $\Bar{\tau}_{P^*}$ is the preferred one. Finally, to obtain an efficient perturbation time interval, plug $\Psi_T(t\leq P^*)$ and $\Bar{\tau}_{P^*}$ into \cref{eq:LR-prediction-rarepert}, and evaluate the speedup for $\max(t_r,\Bar{\tau}_{P^*}) \leq P \leq P^*$. Choose the time interval that leads to the highest speedup. 

\section{Summary and Conclusion}
We developed an approach based on a recent response theory by Keidar et al.~\cite{keidar2024universal} to design useful perturbations for accelerating the training of NNs.  
To that end, we treated the training as a first-passage process to a target test accuracy.
For the case of a CIFAR-10 classifier using the ResNet-18 architecture, we demonstrated that the unperturbed training test accuracy reached a QSS after a relaxation time of 20-30 epochs. We showed that, as a result, we can predict the mean FPT for a wide range of perturbation times from measurements at a single perturbation time. We focused on three perturbations: S\&P, SR and partial SR, but our method is general and can be used for other perturbations.
Lastly, we proposed a strategy to examine different kinds of perturbations and pick the optimal one and an efficient $P$. We showed that this strategy correctly selects S\&P as a better perturbation than partial SR or SR for a CIFAR-10 classifier and identified a useful perturbation time. 
Our work provides a first-passage framework capable of identifying perturbations that result in a speedup when training an NN model up to a target accuracy. It allows for a more rational design of perturbation protocols, based on physical insights. 

\section{Computational details}
The CIFAR-10 dataset contains 50,000 training and 10,000 test RGB images each belonging to one of 10 different classes. The number of images per class across the train and test sets is equal, and each image is 32$\times$32 pixels in size. For classification, we used a common modified version of the ResNet-18 architecture, because the original architecture was designed for much larger images. The modification includes reducing the kernel size of the first convolution layer from 7 to 3, the stride from 2 to 1, and the padding from 3 to 1. We also removed the subsequent max pooling layer to maintain everything else similar to the original architecture (as was done in the work of Zaidi et al.~\cite{zaidi2023does}). 

We did not use any data augmentation techniques except for standard data normalization (by the channel-wise mean and variance of the train images) for both train and test datasets. We initialized the weights and biases with the default initialization of Pytorch, i.e., uniform distribution, and used the FFCV library for faster training~\cite{leclerc2023ffcv}. We trained each model with an SGD optimizer with a learning rate of 0.02. Implementation of the perturbations was straightforward according to their definitions above. The test accuracy curves for each perturbative protocol with $P=20$ are presented in supporting information.


\section*{Acknowledgements}
B.H. acknowledges support from the Israel Science Foundation (grants No. 1037/22 and 1312/22), the Pazy Foundation of the IAEC-UPBC (grant No. 415-2023), and Tel Aviv University Center for Artificial Intelligence and Data Science (TAD).This project has received funding from the European Research Council (ERC) under the European Union’s Horizon 2020 research and innovation program (grant agreement No. 947731 to S.R.). S.M. acknowledges support from the Quantum Science and Technology Center of Tel Aviv University. B.H. and S.M. thank Sheheryar Zaidi for sharing his code. The authors thank Yohai Bar Sinai, Tomer Koren, and Rotem Widman for fruitful discussions.

\bibliography{refs}

\begin{appendices}

\setcounter{section}{0}
\renewcommand{\thesection}{SI}
%\renewcommand{\thesubsection}{SI.\arabic{subsection}}
\renewcommand{\thesubsection}{\Alph{subsection}}

\renewcommand{\thetable}{S\arabic{table}}
\setcounter{table}{0}  % Reset table counter to start from 1
\renewcommand{\thefigure}{S\arabic{figure}}
\setcounter{figure}{0}



\section{Supplemental Material}

\subsection{Data Collection}

\subsubsection{Obtaining Observational Studies}

This study adopts the definition of {\em an observational study} provided by
the NIH National Library of Medicine\footnote{https://www.ncbi.nlm.nih.gov/mesh/68064888},
which states:
``[...] a work that reports on the results of a clinical study
in which participants may receive diagnostic, therapeutic, or other types of interventions,
but the investigator does not assign participants to specific interventions (as in an interventional study).''

We initially retrieved 160,256 papers from the PubMed database using the query
{\em Observational Study[Publication Type]} via the Entrez system. The majority
of these papers were published in 2014 or later, following the
introduction of the MeSH term {\em Observational Study}. We then excluded
publications that were also labeled as {\em Randomized Controlled Trial} or {\em
Clinical Trial}, leaving us with 155,933 exclusive observational studies. Next,
we removed papers that lacked either a structured abstract with a conclusion
subsection or an English full text, resulting in 104,508 observational studies
with a {\em Conclusion(s)} subsection and English full text. Finally, we
excluded articles whose conclusions did not contain any causal or correlational
claims—typically those focused solely on descriptive findings, recommendations,
study implications, or future work. This filtering process resulted in a final
dataset of 81,132 observational studies for analysis (see
Fig.~\ref{fig__flowchart_data_collection}).


\begin{figure}[!ht]
\centering
	
\begin{tikzpicture}
    \tikzstyle{block} = [rectangle, rounded corners, draw, text width=28em, text centered, minimum height=3.5em]

    % Place nodes
    \node [block] (block1) {160,256 observational studies retrieved from the Entrez pubmed database
    using query {\tt Observational Study[Publication Type]} \\
    (up to 2024-08-26)};
    
    \node [block, below of=block1, yshift=-1.8cm] (block2) {
        155,933 exclusive observational studies
    };
    
    \node [block, below of=block2, yshift=-1.5cm] (block2_3) {
        118,713 have a structured abstract \\ 
    };
    \node [block, below of=block2_3, yshift=-1.5cm] (block3) {
        110,677 papers include a subsection of {\em Conclusion(s)} \\ in their {\em structured abstracts}
    };
    \node [block, below of=block3, yshift=-1.5cm] (block4) {
        104,508 are written in English in their full-text
    };
    \node [block, below of=block4, yshift=-1.5cm] (block5) {
        81,132 papers used in our regression analysis
    };

    \draw[->] (block1) -- (block2) node[anchor=west, midway, align=left] {
        Exclude papers that are labeled as \\ Randomized Controlled Trial or Clinical Trial
    };
    
    \draw[->] (block2) -- (block2_3) node[anchor=west, midway, align=left] {
        Exclude papers that do not have structured abstract 
    };

    \draw[->] (block2_3) -- (block3) node[anchor=west, midway, align=left] {
        Exclude papers that do not have a \\ conclusion subsection in the abstract
    };

    \draw[->] (block3) -- (block4) node[anchor=west, midway, align=left] {
        Exclude papers that are not written in English \\ in their full-text
    };

    \draw[->] (block4) -- (block5) node[anchor=west, midway, align=left] {
        Exclude papers that do not contain any causal or \\
        associational claims in the conclusion 
    };

\end{tikzpicture}

	\caption{The process of collecting the observational studies for our analysis.}
	\label{fig__flowchart_data_collection}
\end{figure}


\subsubsection{Journal Rank}

We downloaded publicly accessible SciMago Journal Rank (SJR) data for 2023 and
prior years\footnote{https://www.scimagojr.com/journalrank.php}. For journals in
our dataset that could not be matched to the 2023 data via their ISSN, we
attempted to retrieve their information from previous years. If a journal
remained unmatched (which occurred for only 0.35\% of the papers), we assigned
it an SJR score of 0.47, corresponding to the bottom quartile in our list of
2,629 journals, based on the assumption that journals not listed in the SJR
database are likely to have a low rank.

\subsubsection{Author Writing Experience and Name Disambiguation}

An author’s writing experience is measured by the number of observational
studies they are associated with in our dataset.
To disambiguate author names, we used the
Semantic Scholar API\footnote{https://api.semanticscholar.org/} to obtain a
unique ID for each author per article. Author IDs were successfully retrieved
for 99.8\% of the authors; for the remainder, a unique random ID was generated
and assigned.



%%%%%%%%%%%%%%%%%%%%%%%
% 222
\subsection{Data Preprocessing}

\subsubsection{Causal and Correlational Claim Identification} \label{si_causal_or_correlational}

Building on our previous work  \citep{yu2019EMNLPCausalLanguage} 
and the taxonomy introduced in \citep{sumner2014association,li2017nlp}, we
curated a corpus of 3,061 sentences exemplifying causal language. Using this
dataset, we developed a BioBERT-based
model\footnote{https://github.com/junwang4/causal-language-use-in-science} to
classify a sentence in the conclusion subsection of an abstract into four levels
of claim strength: Causal, Conditional Causal, Correlational, or None. 
A “causal” claim uses terms indicating direct causation (e.g., {\em increase,
make, lead to, effective in, contribute to}) or uses the modal verb {\em can}
followed by a causal verb. A “conditional” claim uses qualifiers (e.g., {\em
may, might, appear to, probably}) to tone down the level of certainty, whereas a
“correlational” claim employs language suggesting an association (e.g., {\em
associated with, predictor, linked with}).

For this study, we merged the “conditional causal” and “correlational”
categories into a single “correlational” category, based on previous research
indicating that general readers often have difficulty distinguishing between these
two \citep{adams2017readers, bratton2020causal}. With this modification, we
fine-tuned the BioBERT-based model using the same training corpus and evaluation
procedure described in \citep{yu2019EMNLPCausalLanguage}, achieving a macro-F1
score of 0.89. Moreover, two empirical comparisons with the ChatGPT family of
models demonstrated that our specifically trained model still have advantages over ChatGPT for
this task \citep{kim2023can, chen2024evaluating}.

In our analysis, if a structured abstract’s conclusion subsection contains two
or more sentences (which occurs in 70.2\% of our dataset) and at least one
sentence is classified as causal, the entire conclusion is considered to include
a causal claim. Based on this criterion, 31.0\% of the 81,132 conclusions
contain a causal claim.

%---------------------
% 22b
\subsubsection{Country Label} \label{si_country}

Author affiliations in the PubMed metadata provide a basis for inferring country
information. However, this can be challenging when the country is not explicitly
mentioned (e.g., “Beth Israel Deaconess Medical Center”). To address this, we
developed an algorithm to infer the country associated with an author’s
affiliation, trained on a dataset of 2.7 million organizations with known
country information from the ORCID
database\footnote{https://support.orcid.org/hc/en-us/articles/360006897394-How-do-I-get-the-public-data-file}.

We evaluated the algorithm using 1,000 randomly sampled affiliations from the
PubMed metadata. These affiliations fell into two categories: (1) those with
insufficient information to infer the country (e.g., “Division of Infectious
Diseases”) and (2) those with sufficient information. For the few cases (7
affiliations) in Category 1, the algorithm output low confidence (less than
0.8), whereas for Category 2, all inferences were accurate. For additional
details, please visit our GitHub
page\footnote{https://github.com/junwang4/affiliation-to-country-inference}.

For each paper, if all author affiliations are classified under the same country
and the average confidence score exceeds 0.8, the paper is labeled with that
country; otherwise, it is labeled as “others.”
Fig.~\ref{fig__country_distribution} shows the distribution of countries that
have published at least 300 observational studies.

\begin{figure}[!ht]
    \includegraphics[width=.9\textwidth]{./figure/country_distri_R.pdf}
        \caption{
        Number of observational studies by authors from 27 exclusive countries that
        have published at least 300 papers in our dataset, along with the “others” category.
        The “others” category includes 23,153 observational studies: 
        15,906 from multiple countries,
         4,931 from various other countries 
         (each with fewer than 300 studies published),
          1,539 lacking affiliation data,
           and 777 with low-confidence inference (i.e., confidence less than 0.8).
        }
        \label{fig__country_distribution}
    \end{figure}


%---------------------
% 22c
\subsubsection{Gender Label} \label{si_gender}

Gender is inferred from an author's forename. Using a dataset of 6 million
(forename, gender) pairs from WikiData, we developed a Gradient Boosting-based
algorithm that leverages the last eight letters of forenames as features to
predict the likelihood that a name is male or female.

Our algorithm\footnote{https://github.com/junwang4/name-to-gender-inference}
performs comparably to the best of the five name-to-gender inference tools when
evaluated on a benchmark dataset of 5,779 manually labeled names
\citep{Santamara2018ComparisonAB}. In our implementation, a name is labeled as
male if its male confidence exceeds 0.83 and as female if its female confidence
exceeds 0.71; otherwise, it is classified as gender-unknown. These thresholds
achieve balanced recall (0.86 for both male and female) and high precision (0.98
for male and 0.97 for female). Note that the lower threshold for females
reflects the fact that 75\% of the training data comprises male names.

Fig.~\ref{fig__gender_distribution} illustrates the gender distribution among
first and last authors. The pattern, where men outnumber women in both
roles, aligns with prior reports of gender imbalances in biomedical research
publications \citep{ioannidis2023gender}, 
with a more noticeable difference observed among last authors (about 2.5 men per woman).


\begin{figure}[!ht]
    \centering
    \begin{subfigure}{\textwidth}
        \centering
        \includegraphics[width=0.45\textwidth]{./figure/gender_distribution.pdf}
    \caption{Gender distribution by author position. F: female; M: male; I:
    initials only; L: low-confidence inference. For first authors, 37.4\% of
    non-initial forenames are identified as female and 52.3\% as male; for last
    authors, 25.1\% are female and 65.6\% are male. Cases with initials only or
    low-confidence inferences are combined as gender-unknown, comprising 16.9\%
    of first authors and 16.2\% of last authors.}
        \label{fig__gender_distribution_1}
    \end{subfigure}
    \hfill
    \begin{subfigure}{\textwidth}
        \centering
        \includegraphics[width=\textwidth]{./figure/country_gender_firstlast_distri_R.pdf}  % Replace with your image path
	\caption{
      Gender distribution by country (after combining genders of first and last
      authors). F: female; M: male; I: initials only; L: low-confidence
      inference. For mainland China (CN), South Korea (KR), and Taiwan (TW),
      approximately half of the names are difficult to infer from the forename
      alone.
      }
        \label{fig__gender_distribution_2}
    \end{subfigure}

    \caption{Gender distribution}
    \label{fig__gender_distribution}
\end{figure}


\newpage

%%%%%%%%%%%%%%%%%%%%%%%
% 333
\subsection{Logistic Linear Mixed-Effects Regression Analysis}

%\newcommand{\ttt}{\texttt\footnotesize}

\begin{table}[ht]
    \centering
%\Large
    \begin{tabular}{@{}ll@{ }ll@{}}
        %\multicolumn{4}{l}{\large Does\_The\_Observational\_Study\_Use\_Causal\_Claim\_In\_Its\_Abstract\_Conclusion$\sim$} \\ [.1in]
        %\multicolumn{4}{l}{\large Has\_causal\_claim\_in\_abstract\_conclusion $\sim$} \\ [.2in]
        \multicolumn{4}{l}{\large Has\_causal\_language\_use\_in\_the\_conclusion\_of\_a\_structured\_abstract $\sim$} \\ [.2in]

       % &&  $\sum_i StudyCharacteristics_i$ &  {\ttt// $i\in$\{Observational, Comparative, Multicenter, Systematic Review, Validation, Evaluation, Case Reports\}} \\

%& &   $\sum_i StudyDesign_i$ & {\ttt // $i\in$\{Retrospective, Prospective, Case-Control, Cross-Sectional, Cohort, Follow-Up, Longitudinal\}}\\

        & \multicolumn{3}{l}{\ttt \large // VARIABLES OF INTEREST} \\ [.1in]
        &  & Num\_coauthors &   {\ttt // log-scaled}  \\
        & + & Num\_obs\_studies\_by\_{\bf first}\_author &  {\ttt// log-scaled} \\
        & + & Num\_obs\_studies\_by\_{\bf last}\_author &  {\ttt// log-scaled} \\
        & + & Gender\_of\_{\bf first}\_author & {\ttt // categorical: F, M, Unknown} \\
        & + & Gender\_of\_{\bf last}\_author & {\ttt // categorical: F, M, Unknown} \\
        & + & Country\_of\_authors & {\ttt// categorical: 27 exclusive countries and others}\\ [.2in]

    & \multicolumn{3}{l}{\ttt \large // CONTROL VARIABLES} \\ [.1in]
        & + & Journal\_rank &   {\ttt// log-scaled: Scimago journal rank}  \\
        & + & Num\_obs\_studies\_in\_the\_journal &   {\ttt// log-scaled}  \\
        & +& $\sum_{i=1}^{N} \textrm{MeSH\_term}_i$  & {\ttt // binary: N$=$771 MeSH terms, each with at least 150 papers} \\[.2in]

    & \multicolumn{3}{l}{\ttt \large // RANDOM EFFECT VARIABLES} \\ [.1in]
& + & (1 $|$ Year) & {\ttt // categorical: the majority are from 2013 to 2024} \\
& + & (1 $|$ Journal) & {\ttt// categorical: 2629 journals} \\  
& + & (1 $|$ Num\_sentences\_in\_the\_conclusion) &  {\ttt// categorical: the majority are 1, 2, or 3 sentences}
\\ [.2in]
\end{tabular}
\caption{
The formula for the logistic linear mixed-effects regression model.
}
\label{tab_modelformula}
\end{table}



The results presented in Table~\ref{table__overall_effects} are based on a
logistic linear mixed-effects regression model (formula provided in
Table~\ref{tab_modelformula}). Specifically, we used the $glmer()$ function from
the R package $lme4$ \citep{Bates2014FittingLM} on 81,132 observations. In our
model, the dependent variable indicates whether the conclusion subsection of an
article's structured abstract contains at least one sentence with a causal
claim. The primary explanatory variables include: (1) Author writing experience
(measured as the number of observational studies published by the first/last
authors, log-transformed to address skewness); (2) Author gender (for first and
last authors); (3) Author country; and (4) Team size (the number of authors,
also log-transformed).



\subsection{Effects of Various Research Topics, Methods, and Concepts} \label{meshtermeffect}

Our dataset contains over 771 Medical Subject Headings (MeSH) terms, each
treated as a control variable, with its absence serving as the reference
category. A MeSH term is included if it is assigned to at least 150 articles in
the dataset and is not a country name (since author country is analyzed
separately). Table~\ref{meshtermeffect} reports the effects of these MeSH terms,
ordered by effect size.

Refer to the following GitHub repository for a detailed report on the effects of
these MeSH terms.

\noindent\href{https://github.com/junwang4/causal-language-use-80k-observational-studies/blob/main/pdf/meshterm_effects.pdf}{https://github.com/junwang4/causal-language-use-80k-observational-studies}

\end{appendices}

\end{document}
